\uuid{wcaE}
\exo7id{1277}
\titre{exo7 1277}
\auteur{legall}
\organisation{exo7}
\datecreate{1998-09-01}
\isIndication{false}
\isCorrection{false}
\chapitre{Calcul d'intégrales}
\sousChapitre{Autre}
\cpge{mpsi-mp2i.derivabilite.c, pcsi.derivabilite.c, ptsi.derivabilite.c, mpsi-mp2i.analyse-asymptotique.b, pcsi.analyse-asymptotique.b, ptsi.analyse-asymptotique.b, mpsi-mp2i.integration.c, pcsi.integration.b, ptsi.integration.b, bcpst1.integration-segment.a, mp-mpi.fonctions-vectorielles.d, mp-mpi.fonctions-vectorielles.e}
\module{Analyse}
\niveau{L1}
\difficulte{}

\contenu{
\question{
Soit $  f :[a,b]\rightarrow { \Rr}  $ une fonction de classe
$  C^1  $ telle que $  f(a)=f(b)=0  .$ Posons ${M=\sup _{x\in [a,b]}
\vert f'(x)\vert }  .$ Montrer que
$  \displaystyle{
\vert \int _a^bf(t)dt\vert \leq M\frac{ (b-a)^2}{ 4} }  .$ (Indication~: faire
des d\'e\-ve\-lop\-pements limit\'es de
$  \displaystyle{
x\mapsto  \int _a^xf(t)dt }  $ et $  \displaystyle{
x\mapsto  \int _x^bf(t)dt }  ).$
}
}
